\subsubsection{Pre-processing}
\label{subsubsec:preprocessing}

To get a single representative value for each path, we compute the average path power, delay, angles of arrival, and angles of departure.
For each path $\ell$, the average path power is computed from channel coefficients, average delay from the propagation delay, average angles of arrival are computed.
\begin{equation}
    \overline{P_\ell} = \frac{1}{|I|}\sum_{i\in1}{|a_{i}^{2}}|,
    \qquad
    \overline{\tau_\ell} = \frac{1}{|I|}\sum_{i\in I}{|\tau_{\ell,i}|}
\end{equation}
\begin{equation}
    \overline{\theta_\ell} = \frac{1}{|I|}\sum_{i\in I}{\theta_{\ell,i}},
    \qquad
    \overline{\phi_\ell} = \frac{1}{|I|}\sum_{i\in I}{\phi_{\ell,i}}
\end{equation}

Paths with $\overline{\tau_\ell} \le 10^{-9}$s are dropped to remove paths with almost zero delays to keep only paths that likely traveleda meaningful distance.
Then, the distance or range for each path is computed by converting remaining path delays using the speed of light.
\begin{equation}
    L_\ell = c\overline{\tau_\ell}
\end{equation}

Instead of a fixed limit, we use a power threshold that adjusts based on distance so we do not lose weak signals from far away.
\begin{equation}
    s_\ell = \max\!\left(\frac{L_\ell}{10},1\right)^{\!1.1},
    \qquad
    P_{\mathrm{th},\ell} = \frac{P_{\min}}{s_\ell},
\end{equation}
Paths with $\overline{P_\ell}>P_{\mathrm{th},\ell}$ are kept.
Finally, each remaining angle pair is converted to a 3D unit direction,
\begin{equation}
    \mathbf{d}_\ell=[\sin\theta_\ell\cos\phi_\ell,\sin\theta_\ell\sin\phi_\ell,\cos\theta_\ell]^{\!\top}
\end{equation}